\section{From a constructed controller to an adoption decision}\label{sec:displacement75}
The attained readout objective controls a uniform residual account. It does not, by itself, identify the expected cost of the controller constructed from that readout. We now qualify that cost directly. A finite acquired controller supplies information that a global payoff bound ignores: two controllers with nearby actions cannot separate arbitrarily under a common innovation. This information remains available when both controllers are discontinuous.

\subsection{A computable pathwise range}
Let $\pi,\nu$ be fixed complete controllers, constant on finite observation partitions $\mathcal C_t^\pi,\mathcal C_t^\nu$ of a bounded invariant state domain. Closures are used when computing distances. Their actions are feasible on the entire corresponding cell. Suppose, uniformly over innovations and feasible pairs,
\begin{align*}
 \|f_t(x,a,z)-f_t(y,b,z)\|&\le L_{f,t}\|x-y\|+B_{f,t}\|a-b\|,\\
 |c_t(x,a)-c_t(y,b)|&\le L_{c,t}\|x-y\|+B_{c,t}\|a-b\|,
 \qquad |g(x)-g(y)|\le L_g\|x-y\|.
\end{align*}
All constants concern the original primitives. Write $b_t=\prod_{s<t}\beta_s$. For $r\ge0$, the finite action-displacement modulus is
\begin{equation}\label{eq:cell-modulus75}
 \alpha_t(r)=\max_{\operatorname{dist}(\overline C,\overline D)\le r}
                 \|\pi_t(C)-\nu_t(D)\|,
 \quad (C,D)\in\mathcal C_t^\pi\times\mathcal C_t^\nu.
\end{equation}
It is nondecreasing. Its value at zero need not vanish, even for a discontinuous controller compared with itself: adjacent closed cells can carry different actions. Exact controller identity is checked separately, on a common refinement, rather than inferred from this deliberately conservative maximum.

\begin{theorem}[Finite-controller displacement certificate]\label{thm:displacement75}
Couple the original trajectory laws by a common initial state and common innovations. Set $\delta_0=0$ and, with $D_X$ the state-domain diameter,
\[
 a_t=\alpha_t(\delta_t),\qquad
 \delta_{t+1}=\min\{D_X,L_{f,t}\delta_t+B_{f,t}a_t\}.
\]
If $G^\pi,G^\nu$ denote the two realized discounted costs, then
\begin{equation}\label{eq:displacement-support75}
 |G^\pi-G^\nu|\le K_{\pi\nu}:=
 \sum_{t<T}b_t(L_{c,t}\delta_t+B_{c,t}a_t)+b_TL_g\delta_T.
\end{equation}
Any independently established cost-range bound may reduce $K_{\pi\nu}$. If the complete controllers agree on a common refinement, their coupled costs agree exactly and the bound can be set to zero. No regularity of the acquired actors or access to a policy-value function is required. Direct enumeration uses $O(\sum_t|\mathcal C_t^\pi||\mathcal C_t^\nu|d)$ arithmetic work for rectangular cells in the $\ell^1$ metric; the resulting bound is qualified before cost sampling.
\end{theorem}

\begin{proof}
Suppose the state displacement at date $t$ is at most $\delta_t$. The actual cell pair has closure distance no greater than that displacement, so its action difference is at most $a_t$. The transition inequality and invariance give the next displacement bound. Induction begins at the common initial state. Applying the cost inequalities at every date and summing proves \eqref{eq:displacement-support75}. Identical controllers choose identical actions at identical states, and induction then gives identical trajectories. Rectangle distances are sums of coordinatewise interval distances; enumerating all pairs gives the stated work.
\end{proof}

For the original two-state investment economy, valid constants are
\[
 L_f=\frac{11}{16},\quad B_f=\frac34,\quad
 L_c=4,\quad B_c=\frac{13}{16},\quad L_g=5,\quad D_X=2.
\]
The supplement derives them directly from the bilinear transition, quartic adjustment cost and shortage term. The common innovation cancels only in the displacement inequality; both marginal trajectories retain the original continuous innovation law. The certificate is thus not a replacement of the stochastic economy by its conditional mean. Its finite pair computation is paid, and can remain expensive when observation partitions are large.

\subsection{Numerical width and economic selection}
Condition on the frozen controllers. For each of $M$ declared estimands $\mu=\mathbb E D$, let independent continuous-bin rows produce endpoints $L_i\le D_i\le U_i$ in a prequalified support $[a,b]$. Define $R=b-a$, the unbiased sample variances $s_L^2,s_U^2$, and $\ell\ge\log(4M/\gamma)$. The endpoint empirical-Bernstein interval is
\begin{equation}\label{eq:cost-interval75}
 [l,u]=[\overline L-r_L,\overline U+r_U]\cap[a,b],\qquad
 r_V=\sqrt{2s_V^2\ell/n}+\frac{7R\ell}{3(n-1)}.
\end{equation}
The standard bounded-sample inequality of \citet{maurer2009}, applied to both endpoints and all estimands, gives simultaneous coverage at least $1-\gamma$. Independence is between rows; different controllers may share every row. Before support clipping, the interval width is exactly
\begin{equation}\label{eq:width-ledger75}
 (\overline U-\overline L)
 +\sqrt{2s_L^2\ell/n}+\sqrt{2s_U^2\ell/n}
 +\frac{14R\ell}{3(n-1)}.
\end{equation}
The first term contains bin, acquisition-branching and arithmetic enclosures. Its separate realized contributions can be audited without claiming that a larger sample removes acquisition discontinuities. The exact integer calculation below records both sampling radii, the range term and mean bin width. Support clipping cannot enlarge the interval. A finite-sample interval is not made valid at arbitrary stopping times by inspecting it repeatedly; the predictable construction in Section~\ref{sec:sequential73} has its own assumptions.

For adoption, write $[l_k,u_k]$ for the gain $\mathbb E(G^0-G^k)$. Let $D$ be the intended number of future deployments, $\lambda\ge0$ the value of a unit of normalized loss, $K_k\ge0$ the priced one-off construction, qualification and installation cost, and $d_k$ the priced incremental cost per deployment. On the same coverage event the total net adoption benefit belongs to
\begin{equation}\label{eq:adoption75}
 [D(\lambda l_k-d_k)-K_k,\ D(\lambda u_k-d_k)-K_k].
\end{equation}
Consequently $D>K_k/(\lambda l_k-d_k)$ supports strictly beneficial adoption whenever the denominator is positive. A negative upper endpoint rules it out. This statement is simultaneous in the nonnegative economic prices and deployment counts; choosing them after inspection does not create an uncovered cost estimand. They remain external economic inputs, and the inference sample size $n$ is not the deployment count $D$. The preceding readout, feasibility, support and inference results now terminate in a decision concerning the actual returned NBO controller.
